\chapter*{Glossari català--anglès}
\addcontentsline{toc}{chapter}{Glossari català--anglès}
\markboth{Glossari català--anglès}{Glossari català--anglès}

Aquest glossari recull la terminologia menys estandarditzada del llibre, amb la seva forma anglesa i, quan escau, la notació emprada. La lectura de bibliografia internacional ---gairebé sempre en anglès--- fa útil tenir a mà aquesta correspondència.

\medskip
\noindent\textbf{Criteri d'anglicismes.} Diversos termes s'han manllevat de l'anglès per manca d'una forma catalana establerta o per fidelitat a la bibliografia (\emph{pullback}, \emph{pushout}, \emph{slice}). En la seva primera aparició es dona, quan existeix, el terme català al costat de l'anglès i de la notació; després se n'usa una sola forma. La terminologia general segueix el \emph{Diccionari de la llengua catalana} de l'Institut d'Estudis Catalans \cite{diec}, i la matemàtica, el \emph{Diccionari de matemàtiques i estadística} del TERMCAT \cite{termcat}.

\medskip

\renewcommand{\arraystretch}{1.25}
\noindent\begin{tabularx}{\linewidth}{@{}>{\raggedright\arraybackslash}X >{\raggedright\arraybackslash}X l@{}}
\hline
\textbf{Català} & \textbf{Anglès} & \textbf{Notació}\\
\hline
categoria oposada & opposite category & $\cat{C}\op$\\
categoria prima & thin category & \\
categoria slice (o llesca) & slice category & $\cat{C}/A$\\
categoria coslice & coslice category & $A/\cat{C}$\\
homconjunt & hom-set & $\Hom(A,B)$\\
functor representable & representable functor & p.\,ex.\ $\Hom(-,A)$\\
plançó de Yoneda & Yoneda embedding & \\
quasiinvers (d'una equivalència) & quasi-inverse / pseudo-inverse & \\
prefeix & presheaf & $\cat{Set}^{\cat{C}\op}$\\
producte fibrat (pullback) & pullback & \\
suma amalgamada (pushout) & pushout & \\
con & cone & \\
cocon & cocone & \\
límit / colímit & limit / colimit & $\varprojlim$ / $\varinjlim$\\
unitat / counitat & unit / counit & $\eta$ / $\varepsilon$\\
composició horitzontal & horizontal composition & $\theta\ast\eta$\\
whiskering (composició lateral) per l'esquerra & left whiskering & $K\eta$\\
whiskering (composició lateral) per la dreta & right whiskering & $\eta H$\\
epimorfisme regular & regular epimorphism & \\
parella nucli & kernel pair & \\
factorització imatge & image factorization & \\
reindexació & reindexing / base change & $f^{*}$\\
subobjecte & subobject & $\Sub(A)$\\
classificador de subobjectes & subobject classifier & $\Omega$\\
sedàs & sieve & \\
semireticle implicatiu & implicative meet-semilattice & \\
categoria cartesiana tancada & cartesian closed category & \\
categoria regular & regular category & \\
\hline
\end{tabularx}
\renewcommand{\arraystretch}{1.0}
